% ============================================================================
%  Chapitre 3, partie I : Modèle et configuration
%  (manuscrit Modélisation, p. 1 à 3)
% ============================================================================

% ----------------------------------------------------------------------------
\section{Modèle et configuration}
% ----------------------------------------------------------------------------

% ----------------------------------------------------------------------------
\subsection{Définition de la sphère de Chaplygin}

\begin{hypotheses}[Sphère de Chaplygin]
    La sphère de Chaplygin est un solide de masse $m$ qui vérifie :
    \begin{itemize}
        \item sa surface est une sphère de rayon $a$.
        \item elle roule sans glisser sur un plan horizontal fixe.
        \item son centre de masse coïncide avec son centre géométrique
            $G$.
        \item sa répartition de masse n'est pas uniforme : ses moments
            principaux d'inertie $I_1$, $I_2$, $I_3$ sont différents.
    \end{itemize}
\end{hypotheses}

\begin{figure}[H]
    \centering
    % ============================================================================
%  Figure : sphère de Chaplygin sur le plan horizontal, repère sol et
%  repère corps
% ============================================================================
\begin{tikzpicture}[>=Stealth, line cap=round]
    \def\a{1.5}   % rayon de la sphère

    % sol hachuré
    \fill[pattern=north east lines] (-3.4,-0.25) rectangle (3.6,0);
    \draw[thick] (-3.4,0) -- (3.6,0);

    % repère sol
    \begin{scope}[shift={(-2.9,0.45)}]
        \draw[->, thick] (0,0) -- (0.9,0) node[below] {$\vect{e}_x$};
        \draw[->, thick] (0,0) -- (0.5,0.4) node[right] {$\vect{e}_y$};
        \draw[->, thick] (0,0) -- (0,0.9) node[left] {$\vect{e}_z$};
        \node[below left] at (0,0) {$O$};
    \end{scope}

    % sphère
    \shade[ball color=accent!12] (0,\a) circle (\a);
    \draw[thick] (0,\a) circle (\a);
    \draw[gray, densely dashed] (-\a,\a) arc (180:360:\a cm and 0.35cm);
    \fill (0,\a) circle (1.5pt);
    \node[below left] at (0,\a) {$G$};

    % repère corps
    \draw[->, thick, coulprop] (0,\a) -- ++(25:0.9)
        node[above right=-2pt] {$\vect{b}_1$};
    \draw[->, thick, coulprop] (0,\a) -- ++(140:0.9)
        node[above left] {$\vect{b}_2$};
    \draw[->, thick, coulprop] (0,\a) -- ++(80:1.0)
        node[above] {$\vect{b}_3$};

    % vecteur GP et point de contact
    \draw[->, coulthm] (0.08,\a) -- (0.08,0.06)
        node[pos=0.6, left] {$-a\,\vect{e}_z$};
    \fill (0,0) circle (1.5pt);
    \node[below=2pt] at (0,-0.25) {$P$};

    % vitesse angulaire
    \draw[->, very thick, couldef] (0,\a) -- ++(-25:2.1)
        node[right] {$\vect{\Omega}$};
\end{tikzpicture}

    \caption{Sphère de Chaplygin de centre $G$ et de rayon $a$, en
    contact avec le sol au point $P$. Le repère corps
    $(\vect{b}_1,\vect{b}_2,\vect{b}_3)$ suit les axes principaux
    d'inertie.}
    \label{fig-sphere}
\end{figure}

% ----------------------------------------------------------------------------
\subsection{Repères}

On utilise les deux repères orthonormés directs du chapitre
précédent :
\begin{itemize}
    \item $\rep{s} = (O,\vect{e}_x,\vect{e}_y,\vect{e}_z)$, lié au sol et
        galiléen, avec $\vect{e}_z$ vertical vers le haut ;
    \item $\rep{b} = (G,\vect{b}_1,\vect{b}_2,\vect{b}_3)$, lié à la
        sphère, dont les axes sont les axes principaux d'inertie.
\end{itemize}
L'orientation de la sphère est la matrice
\[
    \mat{R} = \mat{R}_{b\to s}\in\SO,
    \qquad
    \vect{u}_s = \mat{R}\,\vect{u}_b,
    \qquad
    \mat{R}^T = \mat{R}_{s\to b}
\]
Conformément au chapitre 2, $\vect{\omega}$ désigne la vitesse
angulaire de la sphère exprimée dans $\rep{b}$, et
$\vect{\Omega} = \mat{R}\,\vect{\omega}$ la même vitesse exprimée dans
$\rep{s}$. Dans $\rep{b}$, le tenseur d'inertie est constant et
diagonal :
\[
    \mat{I}_b = \mathrm{diag}(I_1,I_2,I_3)
\]

% ----------------------------------------------------------------------------
\subsection{Espace de configuration}

La position de $G$ ou d'un point matériel et
l'orientation $\mat{R}$. Le contact avec le sol impose que le centre
reste à la hauteur $a$ :
\[
    z_G(t) = a \qquad \forall\,t
\]
C'est une contrainte holonome, qui retire une coordonnée nécessaire. Il reste donc
$6 - 1 = 5$ degrés de liberté :
\[
    q = (x,y,\mat{R})\in Q = \underbrace{\R^2}_{\text{translation plane
    de } G}\times\underbrace{\SO}_{\text{orientation}}
\]
La symétrie géométrique de la sphère ne réduit pas $Q$ : comme la
masse est répartie de façon anisotrope, toutes les orientations sont
distinctes. Supposons que sur la sphère, on marque des points de repères en couleurs différentes pour déterminer l'orientation de la sphère.

% ----------------------------------------------------------------------------
\subsection{Angles d'Euler}

Pour déterminer l'orientation de la sphère, on utilise les angles d'Euler
$(\varphi,\theta,\psi)$ selon la convention $Z$-$Y'$-$X''$ :
\begin{itemize}
    \item $\psi$, l'azimut (\emph{yaw}), autour de la verticale
        $\vect{e}_z$.
    \item $\theta$, le tangage (\emph{pitch}), autour du nouvel axe
        $\vect{a}_2$.
    \item $\varphi$, le roulis (\emph{roll}), autour du dernier axe
        $\vect{c}_1 = \vect{b}_1$.
\end{itemize}
Les coordonnées généralisées sont alors
\[
    q = (x,y,\varphi,\theta,\psi)\in\R^5
\]

\subsubsection{Rotations successives}

On passe de la base du sol à celle du corps par trois rotations
élémentaires, chacune autour d'un axe de la base précédente :
\[
    (\vect{e}_x,\vect{e}_y,\vect{e}_z)
    \xrightarrow[\text{autour de } \vect{e}_z]{\psi}
    (\vect{a}_1,\vect{a}_2,\vect{a}_3)
    \xrightarrow[\text{autour de } \vect{a}_2]{\theta}
    (\vect{c}_1,\vect{c}_2,\vect{c}_3)
    \xrightarrow[\text{autour de } \vect{c}_1]{\varphi}
    (\vect{b}_1,\vect{b}_2,\vect{b}_3)
\]
avec $\vect{a}_3 = \vect{e}_z$, $\vect{c}_2 = \vect{a}_2$ et
$\vect{b}_1 = \vect{c}_1$. Les rotations élémentaires sont
\[
    \mat{R}_x(\varphi) = \begin{pmatrix}
        1 & 0 & 0\\
        0 & \cos\varphi & -\sin\varphi\\
        0 & \sin\varphi & \cos\varphi
    \end{pmatrix},
    \quad
    \mat{R}_y(\theta) = \begin{pmatrix}
        \cos\theta & 0 & \sin\theta\\
        0 & 1 & 0\\
        -\sin\theta & 0 & \cos\theta
    \end{pmatrix},
    \quad
    \mat{R}_z(\psi) = \begin{pmatrix}
        \cos\psi & -\sin\psi & 0\\
        \sin\psi & \cos\psi & 0\\
        0 & 0 & 1
    \end{pmatrix}
\]

\subsubsection{Matrice d'orientation}

Selon la convention d'orientation $Z$-$Y'$-$X''$, on obtient la matrice d'orientation 
$\mat{R}$ : 
\[
    \mat{R} = \mat{R}_z(\psi)\,\mat{R}_y(\theta)\,\mat{R}_x(\varphi)
\]
Avec $c_\alpha = \cos\alpha$ et
$s_\alpha = \sin\alpha$, on a :
\[
    \mat{R} = \begin{pmatrix}
        c_\psi c_\theta
            & c_\psi s_\theta s_\varphi - s_\psi c_\varphi
            & c_\psi s_\theta c_\varphi + s_\psi s_\varphi\\
        s_\psi c_\theta
            & s_\psi s_\theta s_\varphi + c_\psi c_\varphi
            & s_\psi s_\theta c_\varphi - c_\psi s_\varphi\\
        -s_\theta & c_\theta s_\varphi & c_\theta c_\varphi
    \end{pmatrix}
\]

\subsubsection{Vitesse angulaire}

Chaque angle tourne autour de son propre axe, donc la vitesse angulaire
est la somme des trois vitesses élémentaires :
\[
    \vect{\Omega} = \dot\psi\,\vect{e}_z + \dot\theta\,\vect{a}_2
    + \dot\varphi\,\vect{c}_1
\]
Dans $\rep{s}$, $\vect{a}_2 = \mat{R}_z(\psi)\,\vect{e}_y$ et
$\vect{c}_1 = \mat{R}_z(\psi)\mat{R}_y(\theta)\,\vect{e}_x$, d'où
\[
    \vect{\Omega}_s = \begin{pmatrix}
        \Omega_x\\ \Omega_y\\ \Omega_z
    \end{pmatrix}
    = \begin{pmatrix}
        c_\psi c_\theta & -s_\psi & 0\\
        s_\psi c_\theta & c_\psi & 0\\
        -s_\theta & 0 & 1
    \end{pmatrix}
    \begin{pmatrix}
        \dot\varphi\\ \dot\theta\\ \dot\psi
    \end{pmatrix}
\]
Dans le corps, $\vect{\omega} = \mat{R}^T\vect{\Omega}$ donne
\[
    \vect{\omega} = \begin{pmatrix}
        \omega_1\\ \omega_2\\ \omega_3
    \end{pmatrix}
    = \begin{pmatrix}
        1 & 0 & -s_\theta\\
        0 & c_\varphi & s_\varphi c_\theta\\
        0 & -s_\varphi & c_\varphi c_\theta
    \end{pmatrix}
    \begin{pmatrix}
        \dot\varphi\\ \dot\theta\\ \dot\psi
    \end{pmatrix}
    \longleftrightarrow
    \left\{
    \begin{aligned}
        \omega_1 &= \dot\varphi - \dot\psi\,s_\theta\\
        \omega_2 &= \dot\theta\,c_\varphi + \dot\psi\,s_\varphi c_\theta\\
        \omega_3 &= -\dot\theta\,s_\varphi + \dot\psi\,c_\varphi c_\theta
    \end{aligned}
    \right.
\]
